% Study sheet: Volume II, Chapter 16, Responsible AI.
% Numbers: \val{...} from responsible_ai_studysheet_numbers.tex, written by responsible_ai_studysheet_numbers.py.
\documentclass{studysheet}
\input{responsible_ai_studysheet_numbers}
\sheetvolume{2}
\sheetchapter{16}{Responsible AI}
\sheetsources{Numbers from mlsysim and the chapter's scenarios (loan, monitoring, unlearning, radiology); assumed inputs marked.}

\begin{document}
% ===================================================================== FRONT: LEARN
\frontside
\begin{bigidea}
An accurate, fast model can be unfit to deploy: its errors may fall on people who cannot
challenge them, and fleet scale multiplies their reach. Responsible AI turns principles into \hl{measured claims}, budgeted \hl{safeguards} and an \hl{owner} who can stop the
system. In C\textsuperscript{3} terms, responsibility is a coordination constraint on when compute may run.
\end{bigidea}

\begin{leftcol}
\sheetsec{Key ideas}
\begin{ideas}
\idea{Accuracy is not approval.}{An assurance case (claim, evidence, assumptions, owner,
  invalidation rule) covers one use; a new corpus, region or user base expires it.}
\idea{Fairness criteria compare different rates.}{With unequal base rates and imperfect prediction they
  cannot all hold, so choosing one is a governance decision.}
\idea{A gap is evidence, not a verdict.}{It needs a denominator, a confidence interval and intersectional
  coverage; significance does not name a cause.}
\idea{Safeguards are budgeted capacity.}{Monitoring, explanation and review cost compute, storage
  or people on one path; where they run decides who is protected.}
\idea{Deletion must reach the weights.}{A privacy guarantee bounds exposure, not influence; cheap
  forgetting needs lineage and shards chosen before training.}
\idea{Deployed models reshape their data,}{and reviewers defer: few overrides can mean accepted
  errors, so measure the share of AI errors reviewers correct.}
\end{ideas}

\vspace{10pt}
\sheetsec{Four criteria, four different rates}
\begin{sheettable}{@{}l L L@{}}
\thead{Criterion} & \thead{Equal across groups} & \thead{From counts}\\ \midrule
Demographic parity & approval rate & (TP+FP) $\div$ all\\
Equal opportunity & approval among qualified (TPR) & TP $\div$ (TP+FN)\\
Equalized odds & TPR and false-positive rate & also FP $\div$ (FP+TN)\\
Predictive parity & repaid among approved (PPV) & TP $\div$ (TP+FP)\\
\end{sheettable}\par\vspace{3pt}
{\small Disparate impact = lower approval rate $\div$ higher; below \val{fourfifths} flags a gap to
investigate when approval is the benefit; a screen, not a legal test.\par}
\end{leftcol}%
\begin{rightcol}
\sheetsec{Formulas you will use}
\fcard{Is the gap more than noise?}%
  {\ProportionTestStatistic = \frac{\FirstGroupSampleApprovalRate - \SecondGroupSampleApprovalRate}
   {\sqrt{\PooledSampleApprovalRate(1-\PooledSampleApprovalRate)\left(\frac{1}{\FirstGroupSampleSize}
   + \frac{1}{\SecondGroupSampleSize}\right)}}}%
  {Approval rates $\FirstGroupSampleApprovalRate$, $\SecondGroupSampleApprovalRate$ on
   $\FirstGroupSampleSize$, $\SecondGroupSampleSize$ decisions; pooled $\PooledSampleApprovalRate$ = all approvals $\div$ all decisions. With equal true rates,
   $|\ProportionTestStatistic|$ passes \val{zcrit} only \val{signifpct}\% of the time.}%
  {Independent decisions. Below \val{zcrit} is not evidence of fairness; precision grows as the square root of the sample.}%
  {\begin{tikzpicture}[sheet]
     \draw[dim] (20,0) -- (230,0); \draw[dim] (20,0) -- (20,50);
     \foreach \v in {0,50,100}{
       \draw[dim] ({20+\v*1.7},0) -- ({20+\v*1.7},-2.5);
       \node[small, anchor=north] at ({20+\v*1.7},-4) {\v};}
     \draw[accent, dashed, line width=0.5pt] (20,{\val{zcrit}*10}) -- (230,{\val{zcrit}*10});
     \node[note, anchor=south east] at (230,{\val{zcrit}*10+2.5}) {critical \val{zcrit}};
     \draw[accent, line width=1.1pt] plot[domain=0:120, samples=60, smooth] ({20+\x*1.7},{10*\val{zcoef}*sqrt(\x)});
     \filldraw[fill=accent, draw=accent] ({20+\val{nbig}*1.7},{\val{z}*10}) circle (2pt);
     \filldraw[fill=white, draw=accent] ({20+\val{ncrit}*1.7},{\val{zcrit}*10}) circle (2pt);
     \node[lab, anchor=south east] at ({16+\val{nbig}*1.7},46) {chapter: $\ProportionTestStatistic = \val{z}$};
     \node[note, anchor=north west] at ({22+\val{ncrit}*1.7},{\val{zcrit}*10-3}) {\val{ncrit} per group};
     \node[note, anchor=west] at (24,47) {$\ProportionTestStatistic$ for \val{rateA} vs \val{rateB}};
     \node[note, anchor=north] at (125,-12) {decisions per group};
   \end{tikzpicture}}

\vspace{14pt}
\sheetsec{What safeguards cost}
\begin{sheettable}{@{}l L@{}}
\thead{Safeguard} & \thead{Cost}\\ \midrule
Fairness log & decisions $\times$ share logged $\times$ record: \val{monrate}/s = \val{mondaym}M/day $\times$ \val{monshare}\% $\times$ \val{monrec} B = \val{monlog} GB/day\\
Alerts & groups $\times$ metrics $\times$ false-alarm rate: \val{mongroups} $\times$ \val{monper} $\times$ \val{monfa}\% = \val{monalerts} false alerts per check if no gap is real\\
Full retrain & \$\val{fullm}M, \val{fulldays} days\\
SISA unlearning & one deletion retrains one of \val{shards} shards: \$\val{shardk}k, \val{shardh} h; a batch pays for every shard it touches\\
Explanation & exact: every feature subset, \val{feats} features $\approx$ \val{subsetsm}M evaluations; sampled or tree methods cost less; asynchronous leaves the request path\\
\end{sheettable}\par\vspace{3pt}
{\small An alert goes to an owner to investigate; alone it proves no harm. More metrics need
stricter per-test thresholds.\par}
\end{rightcol}

% ===================================================================== BACK: PRACTICE
\backside
\begin{answers}
\ans{Finish}{$\sqrt{\val{pooled}\times\val{pooledc}\times\val{invsmall}}$ = \val{sesmall};
  \val{gapAB} $\div$ \val{sesmall} = \val{zsmall}, below \val{zcrit} (and approximate at this size): noise
  cannot be ruled out. Not evidence of fairness either: collect more decisions and keep investigating (z card).}
\ans{1}{(a) \val{qdaym}M/day $\times$ \val{qshare}\% = \val{qrecm}M records $\times$ \val{monrec} B = \val{qlog} GB/day.
  (b) \val{qgroups} $\times$ \val{monper} = \val{qmetrics} metrics $\times$ \val{monfa}\% = \val{qalerts} false alerts per check.
  (c) An owner's investigation, with stricter per-test thresholds for so many metrics (cost table).}
\columnbreak
\ans{2}{(a) \$\val{fullm}M $\div$ \val{qshards} = \$\val{qshardk}k; \val{qdayh} h $\div$ \val{qshards} = \val{qshardh} h.
  (b) Every shard: \val{qshards} $\times$ \$\val{qshardk}k = \$\val{qbatchm}M, a full retrain. Only touching
  fewer shards (clustered requests) lowers it (cost table).}
\ans{3}{No. Weights unchanged, but corpus and population are assumptions of the claim, so the
  invalidation rule fires. Revalidate the affected claims; restrict use meanwhile (assurance-case idea).}
\ans{4}{Not shown. Fewer overrides fit fewer AI errors or more accepted errors. On a labeled sample,
  compare AI error rates and the share of AI errors reviewers correct, before and after (oversight idea).}
\end{answers}

\begin{leftcol}
\sheetsec{Worked example}
\begin{practicebox}\boxtitle{Does the loan model pass, and is the gap real?}
\val{sizea} applicants per group, outcomes known for all:\par\vspace{3pt}
\begin{sheettable}{@{}l c c c c@{}}
\thead{Group} & \thead{TP} & \thead{FP} & \thead{FN} & \thead{TN}\\ \midrule
A & \val{tpa} & \val{fpa} & \val{fna} & \val{tna}\\
B & \val{tpb} & \val{fpb} & \val{fnb} & \val{tnb}\\
\end{sheettable}\par\vspace{2pt}
{\small Positive = approved; TP: approved and repaid; FN: rejected but would repay.\par}
\begin{steps}
\item Approval: \val{appra}/\val{sizea} = \val{rateA}, \val{apprb}/\val{sizeb} = \val{rateB}; disparate impact
  \val{rateB} $\div$ \val{rateA} = \res{\val{di}}, below \val{fourfifths}.
\item TPR: \val{tpa}/\val{qualA} = \val{tprA}, \val{tpb}/\val{qualB} = \val{tprB} (gap \val{tprgap}). FPR:
  \val{fpa}/\val{unqualA} = \val{fprA}, \val{fpb}/\val{unqualB} = \val{fprB}. Both fail.
\item PPV: \val{tpa}/\val{appra} = \val{ppvA}, \val{tpb}/\val{apprb} = \val{ppvB}: unequal the other way.
\item $\PooledSampleApprovalRate = \val{pooled}$; $\sqrt{\val{pooled}\times\val{pooledc}\times\val{inv}} = \val{se}$;
  \val{gapAB} $\div$ \val{se} = \res{\val{z}} $>$ \val{zcrit}.
\end{steps}
\takeaway{It fails every criterion, and the data show unequal approval rates; neither fact names the cause.}
\end{practicebox}

\vspace{10pt}
\sheetsec{Your turn: finish it}
\begin{practicebox}\boxtitle{Same rates, \val{nsmall} applicants per group}
The audit saw \val{rateA} against \val{rateB} on \val{nsmall} decisions per group.\par
\textbf{Given:} pooled rate \val{pooled}; $1/\FirstGroupSampleSize + 1/\SecondGroupSampleSize = \val{invsmall}$.\par
\textbf{You:} find $\ProportionTestStatistic$, decide what the gap shows, and say what the team does next.\par\vspace{30pt}
\end{practicebox}
\end{leftcol}%
\begin{rightcol}
\sheetsec{Common mistakes}
\mistake{``False-positive rates match, so the model is fair.''}%
  {The chapter's second loan model approves \val{mfpr}\% of defaulters in both groups, yet carries a
   \val{mapprgap}-point approval gap and a \val{mtprgap}-point TPR gap. Pick the criterion from the harm.}{}
\mistake{``A radiologist checking every AI read catches its misses.''}%
  {Alone: radiologist \val{human}\%, AI \val{ai}\%. Accepting \val{accept}\% of the AI's misses fails
   \val{aimiss} $\times$ \val{acceptfrac} = \val{fail} of tumors: at most \val{team}\%, \val{gainpts} points above the AI alone.}{}
\mistake{``We can add explanations after launch.''}%
  {Inline approximate SHAP adds \val{shaplo}--\val{shaphi} ms to a \val{shapbase} ms request at \val{shaprate}/s:
   \val{shaptotlo}--\val{shaptothi} ms and \val{capplo}--\val{capphi}$\times$ the capacity. Plan placement before launch.}{}

\vspace{4pt}
\sheetsec{Check yourself}
\begin{qs}
\item \assumed{} \val{qrate} decisions/s, \val{qshare}\% logged at \val{monrec} B, \val{qgroups} groups
  $\times$ \val{monper} metrics at \val{monfa}\% each. (a) Log per day? (b) False alerts per check if no
  gap is real? (c) What must an alert trigger?
\item \assumed{} the chapter's retrain (\$\val{fullm}M, \val{fulldays} days) split into \val{qshards} shards.
  (a) Cost and time to forget one person? (b) A day's \val{ureq} requests touch every shard: cost?
\item A support assistant keeps its approved weights but swaps its retrieval corpus and opens in a new region.
  Does the approval still hold? What decides, and what next?
\item After a new interface, reviewers override the AI half as often. Has the AI improved? What
  measurement tells?
\end{qs}
\end{rightcol}
\end{document}
